\section{Introduction}
\label{sec:intro}

The emergence of large language model (LLM)-driven mobile GUI agents represents a qualitative
shift in mobile automation fraud.%
Frameworks such as AutoGLM~\cite{autoglm2024}, AppAgent~\cite{appagent2024}, and Anthropic's
Computer Use~\cite{anthropiccu2024} enable a remote operator to direct a smartphone by
supplying natural-language goals: the agent observes the screen, reasons about the next
action, and injects touch events via the Android Debug Bridge (ADB) or
\texttt{AccessibilityService}.%
These systems are openly available, cost under \$0.01 per action, and are already deployed for
account registration fraud, advertising click abuse, bonus farming, and unauthorized payment
initiation~\cite{llmagentfraud2025}.%
Unlike classical scripted bots that repeat fixed action sequences, LLM agents are adaptive:
they navigate unexpected UI states, solve simple CAPTCHAs, and can be equipped with
human-motion synthesis to mimic plausible touch trajectories.%

\textit{Kinematic detection is insufficient against LLM agents.}%
The state-of-the-art detection approach, represented by Turing Test on Screen~\cite{ttos2026},
analyzes touch kinematic features---coordinates, velocity, inter-touch duration---to distinguish
LLM agents from human operators.%
This framing asks: \textit{does this interaction look like a human?}%
An LLM agent augmented with a motion-synthesis model can answer ``yes'' without requiring
any physical contact: it generates touch events whose coordinate trajectories and timing
distributions match the statistical fingerprint of real users.%
Turing Test on Screen itself acknowledges that \textit{``achieving high-fidelity humanization of
SensorEvents poses significant technical challenges''} and delineates physical sensor spoofing
as an open problem outside its scope.%
BeCAPTCHA~\cite{becaptcha2021}, a precursor using accelerometer and touchscreen biometrics,
was designed for fixed-script bots predating the LLM agent era and is similarly circumventable
by realistic motion synthesis.%

We observe that the right question is not ``does this interaction \textit{look} like a
human?''\ but rather ``did this interaction \textit{physically originate} from a human?''%
Every real finger contact with a smartphone touchscreen initiates a chain of physical
events that are fundamentally inaccessible to software automation.%
The mechanical impulse of a fingertip striking the glass propagates through the phone chassis
and produces a measurable acceleration transient at the MEMS accelerometer---a brief,
sub-100\,ms pulse detectable by the \textit{zero-permission} linear accelerometer
(\texttt{TYPE\_LINEAR\_ACCELERATION})~\cite{androidSensor}.%
ADB-injected and \texttt{AccessibilityService}-dispatched touch events are software
\texttt{InputEvent} objects; they traverse the Android input dispatch pipeline but never
actuate the touchscreen hardware.%
No finger makes contact, no impulse propagates, and the accelerometer remains silent.%
This \textit{physical causality gap} is a structural impossibility for software automation---
unlike kinematic statistics, it cannot be forged by any combination of timing adjustments
or synthetic trajectory generation.%

We present \textit{PhysCausal}, a challenge-response system that exploits the physical
causality gap to detect LLM GUI agents using only the zero-permission accelerometer.%
PhysCausal operates in two complementary modes.%
In \textit{passive mode}, the system continuously correlates the touch-event timestamp stream
$\tau(t)$ against the linear acceleration signal $\delta(t)$: for each reported touch, it
checks whether an accelerometer impulse is present in the subsequent [0,\,80]\,ms window.%
In \textit{active mode}, the server additionally selects a random challenge time
$t^* \in [\frac{T}{4},\frac{3T}{4}]$ after recording begins, prompting a single tap;
the physical impulse at $t^*$ is then verified.%
Because $t^*$ is chosen after the session starts, pre-recorded impulse streams cannot
contain a correctly timed transient.%

A \textit{Physical Causality Consistency Network} (PCCN) jointly processes two input streams:
the linear acceleration signal and the sparse touch-event indicator series.%
A patch-based Transformer encoder extracts temporal features from $\delta(t)$; a cross-modal
attention module queries the accelerometer stream at each touch-event position, detecting
the presence or absence of a causal impulse response.%
A third stream encodes the distribution of inter-touch intervals, capturing the characteristic
bimodal latency signature of cloud-API-backed LLM agents.%
PCCN is trained in three stages: self-supervised masked pre-training on unlabeled
accelerometer data, supervised contrastive fine-tuning on labeled sessions, and
cross-modal head training on touch--accelerometer pairs.%

\textit{Overlay agent detection.}
We introduce a third-class threat not addressed by prior work: the \textit{overlay agent},
in which an LLM agent injects fraudulent actions (e.g., payment confirmation taps)
\textit{on top of} a session where a legitimate user is physically holding the phone.%
This corresponds to the ``authorized push payment'' fraud scenario in which an attacker
remotely controls actions while the victim believes they are the sole operator.%
PCCN detects overlay agents because the server-issued challenge at $t^*$ will be executed
by the ADB injection pathway, leaving no accelerometer impulse at that specific timestamp
even though the broader session contains genuine human IMU activity.%

We evaluate PhysCausal using \textcolor{red}{XX} participants on \textcolor{red}{XX} device
models (Pixel, Galaxy, OnePlus, Xiaomi families) as the human class, and three LLM agent
frameworks (AutoGLM, AppAgent, custom ADB scripts) under A0--A2 adversary conditions as the
agent class.%
Post-hoc attention visualization confirms that PCCN independently concentrates on the
[0,\,80]\,ms post-touch window without explicit supervision.%

\noindent\textbf{Contributions.}
\begin{itemize}[leftmargin=*,topsep=2pt,itemsep=1pt]
  \item \textbf{Physical causality as a liveness primitive.}
        We identify and formalize the mechanical impulse response of real finger contact as
        a zero-permission, structurally unforgeable signal that software injection cannot
        replicate---closing the gap that kinematic-based detection leaves open.

  \item \textbf{Physical Causality Consistency Network (PCCN).}
        A unified passive/active architecture that (i)~detects causal coupling between
        touch events and accelerometer impulses via cross-modal attention, and
        (ii)~encodes LLM inference latency signatures via a timing distribution stream---
        defeating A0--A2 adversaries including timing-mimicry attacks.

  \item \textbf{Three-class threat model and overlay agent detection.}
        We define and evaluate the novel overlay agent threat, where an LLM agent injects
        fraudulent actions into a live human session, and show that the server-issued $t^*$
        challenge defeats it at \textcolor{red}{>XX\%} detection accuracy.

  \item \textbf{Cross-device evaluation and interpretability.}
        We demonstrate \textcolor{red}{>XX\%} accuracy under a zero-overlap device-model
        split and show via attention visualization that PCCN rediscovers the physical
        impulse window without explicit feature engineering.
\end{itemize}
